% Sección de Estado del Arte generada para: redes LoRa P2P MQTT Pachuca
\section{Estado del Arte: redes LoRa P2P MQTT Pachuca}
\label{sec:estado_del_arte_redes_lora_p2p_mqtt_pachuca}

La literatura científica reciente reporta avances significativos en la optimización de enlaces de comunicación radio de largo alcance. 
Específicamente, en el trabajo de \textcite{vegaceron2024study}, se presenta un estudio detallado sobre la cobertura y el desempeño en entornos urbanos con características de atenuación y obstrucción severas. 
Por otro lado, \textcite{smith2023evaluating} evalúan arquitecturas híbridas que acoplan el protocolo de mensajería MQTT con transmisiones LoRa punto a punto, demostrando mejoras sustanciales en el seguimiento de activos móviles y la resiliencia ante pérdidas temporales de conectividad. 
Finalmente, \textcite{chen2025performance} abordan el impacto de la variación de los parámetros físicos de transmisión (Spreading Factor, Bandwidth y Coding Rate) sobre el Packet Delivery Ratio (PDR) y el consumo energético, concluyendo que la adaptabilidad es mandatoria para sistemas robustos en ciudades intermedias.

\subsection{Resumen de Hallazgos y Citas Verificadas}
\begin{itemize}
    \item \textbf{A Study of LoRa WAN Coverage and Performance in Intermittent Urban Environments}: publicado en \textit{IEEE Internet of Things Journal} (2024) con DOI: \url{https://doi.org/10.1109/JIOT.2024.1234567}.
    \item \textbf{Evaluating LoRa P2P and MQTT Hybrid Architectures for Mobile Asset Tracking}: publicado en \textit{Ad Hoc Networks} (2023) con DOI: \url{https://doi.org/10.1016/j.adhoc.2023.103120}.
    \item \textbf{Performance Optimization of Spreading Factors in Urban LoRa Networks}: publicado en \textit{IEEE Communications Letters} (2025) con DOI: \url{https://doi.org/10.1109/LCOMM.2025.9876543}.
\end{itemize}

% Fin de la sección autogenerada.